\documentclass[aspectratio=169]{beamer}

\usepackage[utf8]{inputenc}
\usepackage[T1]{fontenc}
\usepackage[brazil]{babel}
\usepackage{amsmath, amssymb}
\usepackage{microtype}
\usepackage{csquotes}
\usepackage{natbib}
\usepackage{tikz}
\usepackage{graphicx}

% Colors
\definecolor{mblack}{RGB}{20, 20, 20}
\definecolor{mprimary}{RGB}{83, 152, 190}
\definecolor{maccent}{RGB}{222, 165, 75}
\definecolor{mwhite}{RGB}{255, 255, 255}

% Theme
\usetheme{default}
\usecolortheme{default}
\setbeamercolor{normal text}{fg=mblack, bg=white}
\setbeamercolor{frametitle}{fg=white, bg=mprimary}
\setbeamercolor{title}{fg=white}
\setbeamercolor{author}{fg=white}
\setbeamercolor{date}{fg=white}
\setbeamercolor{itemize item}{fg=maccent}
\setbeamercolor{itemize subitem}{fg=mprimary}
\setbeamercolor{enumerate item}{fg=maccent}
\setbeamercolor{enumerate subitem}{fg=mprimary}
\setbeamercolor{block title}{fg=white, bg=mprimary}
\setbeamercolor{block body}{fg=mblack, bg=mprimary!8}
\setbeamercolor{block title alerted}{fg=white, bg=maccent}
\setbeamercolor{block body alerted}{fg=mblack, bg=maccent!10}
\setbeamercolor{alerted text}{fg=maccent}

\setbeamerfont{frametitle}{series=\bfseries, size=\large}
\setbeamerfont{title}{series=\bfseries, size=\LARGE}
\setbeamerfont{block title}{series=\bfseries}

\setbeamertemplate{navigation symbols}{}

% Footer
\setbeamertemplate{footline}{%
  \leavevmode%
  \hbox{%
    \begin{beamercolorbox}[wd=0.7\paperwidth, ht=2.25ex, dp=1ex, leftskip=0.3cm]{normal text}%
      \tiny\color{mblack!50} BCC502 -- Metodologia Científica para Ciência da Computação
    \end{beamercolorbox}%
    \begin{beamercolorbox}[wd=0.3\paperwidth, ht=2.25ex, dp=1ex, rightskip=0.3cm, right]{normal text}%
      \tiny\color{mblack!50}\insertframenumber{} / \inserttotalframenumber
    \end{beamercolorbox}%
  }%
}

% Frametitle with accent rule
\setbeamertemplate{frametitle}{%
  \nointerlineskip%
  \begin{beamercolorbox}[wd=\paperwidth, ht=0.55cm, dp=0.2cm, leftskip=0.5cm]{frametitle}%
    \usebeamerfont{frametitle}\insertframetitle%
  \end{beamercolorbox}%
  \vspace{-0.1cm}%
  \textcolor{maccent}{\rule{\paperwidth}{2pt}}%
}

% Title page: three-band layout
\setbeamercolor{titlebar top}{bg=mwhite}
\setbeamercolor{titlebar main}{bg=mprimary}
\setbeamercolor{titlebar bottom}{bg=maccent}

\setbeamertemplate{title page}{%
  \vspace{0pt}%
  \begin{beamercolorbox}[wd=\paperwidth, ht=0.55cm, dp=0pt]{titlebar top}\end{beamercolorbox}%
  \begin{beamercolorbox}[wd=\paperwidth, ht=3.8cm, dp=0pt, center]{titlebar main}%
    \vspace{0.6cm}
    {\usebeamerfont{title}\color{white}\inserttitle\par}%
    \vspace{0.35cm}{\color{white}\small\insertauthor\par}%
    \vspace{0.15cm}{\color{white}\footnotesize\insertdate\par}%
    \vspace{0.25cm}
  \end{beamercolorbox}%
  \begin{beamercolorbox}[wd=\paperwidth, ht=1.3cm, dp=0pt]{titlebar bottom}\end{beamercolorbox}%
}

\setlength{\itemsep}{4pt}

% Highlight commands
\newcommand{\impt}[1]{\textcolor{maccent}{\textbf{#1}}}
\newcommand{\ntc}[1]{\textcolor{mprimary}{#1}}

\title{BCC502 -- Metodologia Científica}
\author{Rage Against the Inductivism\footnote{Baseado em \citep{chalmers1993ciencia}}}
\date{\today}

\begin{document}

% ==========================================
\begin{frame}[plain]\titlepage\end{frame}

% ==========================================
\begin{frame}{Sumário}\tableofcontents\end{frame}

% ==========================================
\section{Recapitulando o Indutivismo}
% ==========================================

\begin{frame}{Recapitulando o Indutivismo}
  \begin{enumerate}
    \item A ciência começa com observação
    \pause
    \item A observação fornece uma base segura sobre a qual o conhecimento científico pode ser construído
    \pause
    \item O conhecimento científico é obtido a partir de proposições de observação por \impt{indução}
  \end{enumerate}

  \vspace{0.4cm}
  \begin{block}{Princípio da Indução}
    Se um grande número de \(A\)s foi observado sob uma ampla variedade de condições, e se todos esses \(A\)s observados possuíam sem exceção a propriedade \(B\), então \impt{todos} os \(A\)s possuem a propriedade \(B\).
  \end{block}
\end{frame}

% ==========================================
\section{Justificação da Indução}
% ==========================================

\begin{frame}{Argumento Lógico}
  Se a premissa é verdadeira, a conclusão \textbf{deve} ser verdadeira?

  \vspace{0.3cm}
  \begin{itemize}
    \item Argumentos \ntc{dedutivos} possuem esse caráter
    \pause
    \item Argumentos \ntc{indutivos}, \impt{não}
    \pause
    \item É possível que a conclusão de um argumento indutivo seja falsa embora as premissas sejam verdadeiras --- sem contradição
  \end{itemize}

  \vspace{0.3cm}
  \pause
  \begin{alertblock}{Exemplo: Corvos Pretos}
    Não há contradição em afirmar simultaneamente:
    \begin{itemize}
      \item Todos os corvos observados se revelaram pretos
      \item Nem todos os corvos são pretos
    \end{itemize}
  \end{alertblock}
\end{frame}

\begin{frame}{Argumento da Experiência}
  O princípio da indução \textit{funcionou} historicamente em muitos casos:
  \begin{itemize}
    \item Leis da ótica
    \item Leis do movimento planetário
    \item etc.
  \end{itemize}

  \vspace{0.4cm}
  \pause
  \begin{alertblock}{Circularidade (David Hume, séc. XVIII)}
    O próprio argumento indutivo é empregado para validar a indução:
    \[
    \begin{aligned}
     & \text{A indução foi bem sucedida no caso }x_1   \\
     & \text{A indução foi bem sucedida no caso }x_2    \\
     & \cdots  \\
     & \text{A indução foi bem sucedida no caso }x_n   \\
    \Rightarrow \; & \text{A indução é \impt{sempre} bem sucedida}
    \end{aligned}
    \]
  \end{alertblock}
\end{frame}

% ==========================================
\section{Outros Problemas do Indutivismo}
% ==========================================

\begin{frame}{Vagueza das Condições}
  As condições para aceitar generalizações são definidas de forma vaga:

  \vspace{0.3cm}
  \begin{columns}[T]
    \begin{column}{0.45\textwidth}
      \begin{alertblock}{Perguntas sem resposta}
        \begin{itemize}
          \item ``Grande número'' de observações --- \impt{quantas?}
          \item ``Ampla variedade'' de circunstâncias --- \impt{quais?}
        \end{itemize}
      \end{alertblock}
    \end{column}
    \begin{column}{0.50\textwidth}
      \begin{block}{Exemplo: Ponto de fervura da água}
        Circunstâncias relevantes:
        \begin{itemize}
          \item Pressão atmosférica
          \item Pureza da água
          \item Método de aquecimento
          \item Cor do recipiente (?)
        \end{itemize}
      \end{block}
    \end{column}
  \end{columns}

  \vspace{0.3cm}
  \pause
  \begin{itemize}
    \item A escolha das circunstâncias é baseada em \impt{conhecimento prévio}?
    \item Isso não vem \textbf{antes} da observação?
    \item O indutivismo não se torna \impt{subjetivo}?
  \end{itemize}
\end{frame}

% ==========================================
\section{Recuo para Probabilidade}
% ==========================================

\begin{frame}{A Reformulação Probabilística}
  Diante do problema de Hume, os indutivistas recuam:

  \vspace{0.2cm}
  \[
    \text{Observações repetidas} \;\Rightarrow\; \text{A lei é \textit{provavelmente} verdadeira}
  \]

  \vspace{0.3cm}
  \pause
  \begin{alertblock}{Por que isso não funciona?}
    \begin{enumerate}
      \item Uma afirmação universal se aplica a \impt{infinitos} casos possíveis
      \item O número de observações realizadas é sempre \impt{finito}
      \item Portanto:
      \[
        P = \frac{\text{casos confirmados (finito)}}{\text{casos possíveis (infinito)}} = \frac{n}{\infty} \approx 0
      \]
      \item A probabilidade permanece \impt{próxima de zero}, por mais observações que se acumulem
    \end{enumerate}
  \end{alertblock}
\end{frame}

\begin{frame}{Consequências do Recuo Probabilístico}
  \begin{itemize}
    \item O recuo para a probabilidade \impt{não salva} o indutivismo
    \pause
    \item Substituir ``certeza'' por ``alta probabilidade'' fracassa no quadro lógico-formal
    \pause
    \item Enunciados universais continuam tendo probabilidade essencialmente \impt{nula} com base em qualquer conjunto finito de observações
  \end{itemize}

  \vspace{0.4cm}
  \pause
  \begin{block}{Conclusão}
    Nem a versão forte (``a indução \textbf{prova} leis'') nem a versão fraca (``a indução torna leis \textbf{prováveis}'') se sustentam como justificação lógica do conhecimento científico.
  \end{block}
\end{frame}

% ==========================================
\section{Respostas ao Problema da Indução}
% ==========================================

\begin{frame}{Respostas ao Problema da Indução}
  Duas saídas possíveis:

  \vspace{0.4cm}
  \begin{enumerate}
    \item \textbf{Posição de Hume:} A ciência se baseia na indução, mas a indução não pode ser justificada racionalmente
    \begin{itemize}
      \item A crença em leis científicas não passa disso --- uma \impt{crença}
    \end{itemize}

    \vspace{0.3cm}
    \pause
    \item \textbf{Negar a indução:} A ciência \impt{não} é baseada em indução
    \begin{itemize}
      \item Alternativas: falsificacionismo (Popper), paradigmas (Kuhn), programas de pesquisa (Lakatos)
    \end{itemize}
  \end{enumerate}
\end{frame}

% ==========================================
\section{A Dependência da Observação em relação à Teoria}
% ==========================================

\begin{frame}{Suposições do Indutivismo sobre Observação}
  \begin{block}{Duas suposições centrais}
    \begin{enumerate}
      \item ``A ciência começa na observação''
      \item ``A observação produz uma base segura da qual o conhecimento pode ser derivado''
    \end{enumerate}
  \end{block}

  \vspace{0.3cm}
  \pause
  Premissas implícitas:
  \begin{itemize}
    \item O observador tem acesso mais ou menos direto ao mundo externo pela visão
    \item Dois observadores normais vendo o mesmo objeto ``verão'' a mesma coisa
  \end{itemize}

  \vspace{0.3cm}
  \pause
  \impt{Será que isso é verdade?}
\end{frame}

\begin{frame}{Experiências Visuais $\neq$ Imagens na Retina}
  Há evidência abundante de que dois observadores normais, vendo o mesmo objeto, podem ter experiências visuais \impt{diferentes}.

  \vspace{0.3cm}
  \begin{itemize}
    \item Ilusões de ótica e figuras ambíguas
    \item Radiografias: o que um estudante e um médico ``veem'' é diferente
  \end{itemize}

  \vspace{0.3cm}
  \pause
  \begin{block}{Conclusão}
    A experiência visual:
    \begin{itemize}
      \item \impt{Não} é determinada apenas pelas imagens na retina
      \item Depende da \textbf{experiência prévia} do observador
      \item Depende do \textbf{estado interno} do observador
    \end{itemize}
  \end{block}
\end{frame}

\begin{frame}{Proposições de Observação Pressupõem Teoria}
  \begin{block}{Tese de Chalmers}
    Proposições de observação são \impt{sempre} feitas na linguagem de uma teoria. Não existe observação ``pura''.
  \end{block}

  \vspace{0.3cm}
  \pause
  \textbf{Exemplos:}
  \begin{itemize}
    \item \ntc{Telescópio de Galileu:} pressupunha uma teoria da ótica que garantisse imagens confiáveis
    \pause
    \item \ntc{Estudante de medicina:} um novato e um médico ``veem'' coisas diferentes na mesma radiografia
    \pause
    \item \ntc{Circuitos elétricos:} o novato vê fios; o físico ``vê'' um amplificador
  \end{itemize}
\end{frame}

\begin{frame}{Observações Herdam a Falibilidade das Teorias}
  Se as observações dependem de teoria, herdam sua \impt{falibilidade}.

  \vspace{0.3cm}
  \textbf{Exemplos:}
  \begin{itemize}
    \item \ntc{Torre e a rotação da Terra:} a física aristotélica previa que a pedra ``ficaria para trás'' --- observação baseada em teoria falha
    \pause
    \item \ntc{Varas submersas:} ``a vara está torta'' pressupõe que a visão é infalível, ignorando a refração da luz
    \pause
    \item \ntc{Astronomia pré-telescópica:} planetas ``errantes'' descritos em termos geocêntricos --- reinterpretados com Copérnico
  \end{itemize}

  \vspace{0.3cm}
  \pause
  \begin{alertblock}{Conclusão}
    Toda observação pressupõe um quadro teórico prévio. Não há ponto de partida observacional seguro e independente de teoria.
  \end{alertblock}
\end{frame}

% ==========================================
\section{Indutivismo: Refutado ou Não?}
% ==========================================

\begin{frame}{Indutivismo não Conclusivamente Refutado}
  \begin{itemize}
    \item \textbf{Concessão:} Indutivistas modernos aceitam que a observação depende da teoria
    \pause
    \item \textbf{Manutenção do núcleo:} Para \textit{avaliar} uma teoria, o princípio indutivista permanece: reunir fatos relevantes, sob ampla variedade de condições
    \pause
    \item \textbf{Vulnerabilidade compartilhada:} Alternativas (Popper, Kuhn, Lakatos) também estão sujeitas a críticas sérias
  \end{itemize}

  \vspace{0.4cm}
  \pause
  \begin{block}{Conclusão Prudente}
    O indutivismo ingênuo foi refutado, mas o apelo à evidência empírica acumulada sobrevive como elemento da prática científica. O indutivismo, em versões mais sofisticadas, permanece como posição filosófica legítima --- ainda que \impt{imperfeita}.
  \end{block}
\end{frame}

% ==========================================
\begin{frame}[allowframebreaks]{Referências}
  \bibliographystyle{apalike}
  \bibliography{references}
\end{frame}

\end{document}